% ============================================================
% Script-Aware Bilingual Dyslexia Handwriting Classification
% Draft: ScriptAwareDyslexia_TBD_v1_2026-09
% Venue: TBD. Neutral article class, retargetable.
% Compiler: pdfLaTeX. Citation: numeric natbib (switchable).
% Adapted from academic-paper skill latex_article_template.tex
% (overrides: single spacing, numeric citations, CS structure).
% ============================================================

\documentclass[11pt,a4paper]{article}

\usepackage[margin=1in]{geometry}
\usepackage{amsmath,amssymb}
\usepackage{graphicx}
\usepackage{booktabs}
\usepackage{float}
\usepackage[utf8]{inputenc}
\usepackage[T1]{fontenc}
\usepackage{times}
\usepackage{xcolor}
\usepackage{url}
\usepackage{hyperref}

% --- Citation style: numeric now; switch to author-year by
% --- replacing the two lines below with:
% --- \usepackage[round]{natbib} \bibliographystyle{apalike}
\usepackage[numbers,sort&compress]{natbib}
\bibliographystyle{unsrtnat}

% Visible TODO marker for items the authors must verify or decide.
\newcommand{\todonote}[1]{\textcolor{red}{[TODO: #1]}}

\hypersetup{colorlinks=true,linkcolor=black,citecolor=blue,urlcolor=blue}

\title{Does Script-Aware Conditioning Help? A Controlled Ablation for
Urdu--English Bilingual Dyslexia Handwriting Classification}

\author{%
Ayesha Ashfaq$^{1}$ \and Abdullah Butt$^{1}$\thanks{Corresponding
author.}\\
$^{1}$\todonote{affiliation}%
}

\date{Draft v1, September 2026. Not for distribution.}

\begin{document}
\maketitle

\begin{abstract}
\todonote{Abstract is written last, after results exist. Placeholder
only; no numbers may appear here that do not trace to executed
experiments.}
\end{abstract}

% ============================================================
\section{Introduction}
% ============================================================

Dyslexia is a neurodevelopmental disorder marked by persistent
difficulty with reading, writing, and spelling, and children with
dyslexia frequently exhibit handwriting anomalies such as irregular
spacing, letter reversals, and inconsistent character sizing
\citep{kashif2026}. Because handwriting can be collected with paper,
a pen, and a phone camera, handwriting images are an unusually
accessible signal for early screening support in settings where
neuroimaging or specialist assessment is out of reach
\citep{kashif2026, kunhoth2022}. A growing body of work applies deep
learning to handwriting images for dyslexia and dysgraphia detection
across writing systems, including Latin-script letter corpora
\citep{isa2019, alqahtani2023, robaa2024}, Hindi word samples
\citep{yogarajah2020}, and Chinese dictation responses
\citep{liu2024dysditect}.

Nearly all of this work is monolingual, and the available datasets
and models concentrate on single-script settings. For Urdu, a
low-resource language written in a cursive, dot-bearing Nastaliq
style, the first bilingual children's handwriting dataset for
dyslexia detection appeared only recently: \citet{kashif2026}
collected 852 character-level images (426 from dyslexic and 426 from
non-dyslexic children aged 6--10 in Pakistan), spanning the full
English alphabet, the full Urdu alphabet, and digits, and benchmarked
six script-agnostic convolutional and transfer-learning backbones
under a repeated-run protocol, reporting a best mean accuracy of
$0.7871 \pm 0.0102$ for MobileNetV1 with RMSProp. Their evaluation
treats every character identically regardless of script: one shared
backbone processes Urdu and English characters alike.

This paper asks a deliberately narrow question about that design
choice. Urdu and English character images have sharply different
visual statistics. Urdu characters are cursive, carry
identity-bearing dot (nuqta) configurations, and vary with ligature
context, while English characters are discrete Latin glyphs;
\citet{kashif2026} themselves note that the complexity of Urdu
handwriting raises intra-class variability and complicates feature
extraction. A script-agnostic backbone must devote shared capacity to
both distributions at once. We therefore test whether
\emph{explicitly conditioning the model on the script of each
character image} improves binary dyslexic versus non-dyslexic
classification relative to a script-agnostic baseline of matched
capacity, holding the dataset, the split, and the training protocol
fixed. The script-statistics argument predicts that any benefit
should concentrate in the Urdu subset, so we treat the Urdu-subset
effect as a first-class result alongside the pooled effect rather
than allowing a script-specific gain to be averaged away. The
dataset also contains handwritten digits; we retain digits as a third
conditioning category throughout and defer interpretation of digit
results until the numeral forms present in the images are verified
\todonote{verify from the images whether digits are Western Arabic
(0--9) or Eastern Arabic-Indic glyph forms before interpreting the
digit category}.

The claim we defend is relative, not absolute. The dataset's dyslexic
and non-dyslexic samples were collected from different institutions,
the released images carry no writer identifiers, and consequently
evaluation must use an image-level split. Absolute accuracy under
these conditions reflects within-dataset separability and cannot
support a generalization or screening claim, and we make no such
claim. What the design does support is a controlled comparison: every
experimental arm trains and tests on the identical split under the
identical confounds, so the systematic difference between an arm with
script conditioning and its matched-capacity control is informative
about the conditioning mechanism itself, even though the absolute
numbers are not informative about deployment. We state the boundary
of this logic explicitly in Section~\ref{sec:limitations}: it assumes
the confounds are equally exploitable by both arms, an assumption we
argue is reasonable but cannot prove. We report the matched-pair
difference with repeated seeds, paired statistics, and effect sizes,
and we report it whether it is positive, null, or negative.

The contributions are:
\begin{enumerate}
  \item \textbf{Primary.} A script-conditioned character-level
  classifier for Urdu--English bilingual dyslexia handwriting, and a
  controlled ablation on the dataset of \citet{kashif2026} isolating
  the effect of script conditioning against a matched-capacity
  script-agnostic control under an identical image-level split, with
  both arms tuned with equal effort, conditioning depth swept rather
  than fixed arbitrarily, and enough repeated seeds that a small true
  effect is not lost to run-to-run variance.
  \item \textbf{Supporting.} A re-run of the backbones of
  \citet{kashif2026} under our single fixed split, making all
  comparisons in this paper internally consistent rather than
  borrowed across differing random partitions.
  \item \textbf{Supporting.} A per-script analysis (Urdu, English,
  digits) of every arm, locating where conditioning helps and testing
  the prediction that gains concentrate in Urdu.
\end{enumerate}

The remainder of the paper is organized as follows.
Section~\ref{sec:related} positions this work and, in particular,
disclaims mechanism-level novelty for script-aware computation, which
is established in multi-script recognition. Section~\ref{sec:data}
describes the dataset as released. Section~\ref{sec:method} specifies
the conditioning mechanism and the matched-capacity control.
Section~\ref{sec:setup} fixes the evaluation protocol.
Section~\ref{sec:results} reports results,
Section~\ref{sec:limitations} states limitations, and
Section~\ref{sec:future} lists deferred extensions.

% ============================================================
\section{Related Work}
\label{sec:related}
% ============================================================

\subsection{Handwriting-based dyslexia and dysgraphia detection}

Automated analysis of children's handwriting for learning-disorder
screening spans dysgraphia and dyslexia and a range of capture
modalities; \citet{kunhoth2022} survey approaches from tablet-based
kinematic features to offline handwriting images. For offline images
in Latin script, a widely used line of work builds on a combined
letter corpus assembled from NIST Special Database 19 uppercase
letters, a Kaggle lowercase set, and samples from dyslexic
schoolchildren in Penang, Malaysia, labeled at the letter level as
Normal, Reversal, or Corrected \citep{isa2019, alqahtani2023,
robaa2024}. On this corpus, \citet{isa2019} compared compact CNNs,
\citet{alqahtani2023} reported hybrid CNN--SVM and CNN--RF pipelines,
and \citet{robaa2024} combined transfer learning with Grad-CAM
visualizations, reporting a test precision of 99.65\% and framing
explainability as central to adoption. A recent extension synthesizes
word images from the same letters and applies YOLO-based detection to
localize Normal, Reversal, and Corrected letters within words
\citep{yolo2025}. Beyond Latin script, \citet{yogarajah2020} classified
Hindi word images from children identified by a screening instrument.
Two properties of this literature matter for the present paper.
First, the letter-orientation labels of the Penang-derived corpus are
a different construct from a child-level dyslexia label: they mark
what a letter looks like, not who wrote it. Second, these systems are
monolingual; each model sees one script.

\subsection{Bilingual dyslexia handwriting classification}

The closest prior work is \citet{kashif2026}, who introduced the
first bilingual Urdu--English children's handwriting dataset for
dyslexia detection and evaluated CNN, VGG16, InceptionV3, and three
MobileNet generations, each under three optimizers, with a
repeated-run protocol reporting means, standard deviations, and
confidence intervals. Every backbone in their study is
script-agnostic: Urdu characters, English characters, and digits pass
through one shared network with no script information. Their split is
an 80:10:10 partition at the image level, and they identify the small
dataset size and the absence of participant-wise validation as
limitations of the framework's practical use. We adopt their dataset
as-is, re-run their backbones under our own fixed split for internal
consistency, and add precisely one design factor, script
conditioning, under matched capacity.

\subsection{Script-aware computation in handwriting recognition}

Script-aware processing is not a new mechanism. Multi-script
handwriting recognition has long combined script identification with
recognition, either as cascaded stages or jointly; for example,
MuLTReNets couples script identification and handwriting recognition
in one multilingual network \citep{multrenets2020}. We therefore
claim no novelty for the idea of conditioning computation on script.
The contribution here is of a different kind: an application of
script conditioning to bilingual dyslexia classification, where it
has not previously been evaluated, together with a controlled
ablation that isolates whether it helps in this setting at all. In
recognition, script identity determines the label space, so
script-awareness is structurally necessary; in binary dyslexia
classification the label space is shared across scripts, and whether
script conditioning helps is an open empirical question rather than a
structural requirement. That question is what this paper answers.

\subsection{Attention and conditioning in dyslexia models}

Attention mechanisms have entered dyslexia detection, but not as
script conditioning. DysDiTect \citep{liu2024dysditect} combines a
ResNet-50 feature extractor with positional encoding, a
bidirectional LSTM, and multi-head self-attention over the temporal
sequence of a Chinese dictation task, trained on roughly 100{,}000
character images from 1{,}064 children and reaching 83.2\% accuracy
from handwriting features alone; notably, its data allowed a
participant-level stratified split, which the dataset used here does
not \todonote{keep this contrast; it motivates the Limitations
section}. \citet{robaa2024} use attention-era architectures for
explainability rather than conditioning. In both cases the setting is
a single script, and attention serves sequence modeling or
post-hoc explanation, not script-conditioned feature extraction.

To the best of our knowledge, and based on literature searches
conducted through September 2026 \todonote{authors: re-run this
search at submission time and record queries and dates in the
supplement}, no prior work evaluates script-conditioned computation
for dyslexia classification on real bilingual handwriting. The
nearest neighbors are \citet{kashif2026}, which is bilingual but
script-agnostic, and \citet{liu2024dysditect}, which uses attention
but is monolingual. This paper occupies exactly the cell those two
works leave open, and claims only that cell.

% ============================================================
\section{Data}
\label{sec:data}

We use the bilingual children's handwriting dataset of
\citet{kashif2026} exactly as released, with no exclusions and no
re-collection. The release contains 852 character-level photographs,
organized into two folders: 426 images from children labeled dyslexic
(YES) and 426 from children labeled non-dyslexic (NO). Per the source
study, each participating child, aged 6--10, wrote all letters of the
English alphabet, all letters of the Urdu alphabet, and the digits
0--9 on paper; characters were photographed individually with an
iPhone~11 camera, screened for sharpness and completeness, and, for
Urdu, checked for the visibility of dots, curves, loops, and stroke
formation without any OCR-based correction that could erase
dyslexia-relevant cues \citep{kashif2026}.

Two properties of the release shape everything downstream and are
stated here plainly rather than in fine print. First, the dyslexic
and non-dyslexic classes were sourced from different institutions:
dyslexic samples from a government institute for slow learners in
Narowal and non-dyslexic samples from mainstream schools, with
participants selected via school records and psychological testing
\citep{kashif2026}. The class label is therefore confounded with
collection site, and the label construct is institutional placement
supported by psychological assessment, not an independent clinical
dyslexia diagnosis. Second, the release contains no writer
identifiers and no per-writer grouping; filenames do not link images
to children. Writer-disjoint evaluation is consequently impossible on
this release, which is why our evaluation is defined, and claimed,
strictly within-dataset (Section~\ref{sec:setup}) and why the
comparison logic of this paper is relative rather than absolute
(Section~\ref{sec:method}).

\paragraph{Script annotation.} The release provides the diagnosis
label only; the script of each image is not recorded in metadata. We
therefore annotate every image with a script tag
$s \in \{\textsc{urdu}, \textsc{english}, \textsc{digit}\}$ by visual
inspection. Both authors annotate all 852 images independently;
inter-annotator agreement is reported, and disagreements are resolved
jointly with the resolution logged. The annotation protocol also
records, for every \textsc{digit} image, whether the glyph is a
Western Arabic numeral (0--9) or an Eastern Arabic-Indic numeral
form, since Urdu writing commonly uses the latter; digit-category
results are interpreted only in light of this recorded composition.
\todonote{Run the annotation pass and insert the agreement statistic
and the digit glyph composition here.} Characters whose glyph form is
ambiguous between scripts are flagged at annotation time and reported
as a count. The completed script annotations will be released
alongside our code as a supplement to the original dataset.

\paragraph{Access and ethics.} The source study reports guardian
consent, child assent, institutional approval, and consent to
publication of data in an open-access format \citep{kashif2026}. The
dataset is reachable via the link published in the source paper; we
note that the paper's body describes the data as available at a
public link while its formal Data Availability Statement describes
availability on request, and we treat the published link as the
operative access route. \todonote{Record the access date, retain a
copy of the folder listing as accessed, and, recommended, obtain
brief written confirmation from the corresponding author that use for
this study is within the intended terms.}

% ============================================================
\section{Method}
\label{sec:method}

\subsection{Notation and design principle}

Each sample is a character image $x$ with a binary label
$y \in \{0,1\}$ (dyslexic, non-dyslexic) and an annotated script tag
$s \in \{\textsc{urdu}, \textsc{english}, \textsc{digit}\}$
(Section~\ref{sec:data}). All experimental arms share one design
principle: every arm differs from its comparator by exactly one
factor, and the primary comparison holds parameter count fixed while
varying only the \emph{information} available to the model. Given 852
training-eligible images, all arms are deliberately parameter-light:
a frozen ImageNet-pretrained backbone with small trainable modules,
following the frozen-extractor regime that \citet{kashif2026} found
workable at this scale.

\subsection{Backbone}

The shared backbone is MobileNetV1 pretrained on ImageNet with
convolutional layers frozen, matching the best-performing family in
\citet{kashif2026} so that differences against their re-run setup are
attributable to the added modules rather than to a backbone swap.
Images are resized and normalized identically across all arms
\todonote{fix exact input resolution and normalization once the
re-run pipeline is implemented; mirror the source study where
sensible}. The backbone exposes intermediate feature maps at three
depths, referred to as early, mid, and late block boundaries, where
conditioning modules may be inserted.

\subsection{Script-conditioned attention (SCA) module}

The conditioning mechanism is intentionally simple and
parameter-light. Let $F \in \mathbb{R}^{H \times W \times C}$ be the
feature map at an insertion depth, and let
$e_s \in \mathbb{R}^{d}$ be a learned embedding of the script tag.
The module computes channel-attention gates conditioned on script in
the squeeze-and-excitation style \citep{hu2018squeeze}: the feature
map is globally average-pooled to $z \in \mathbb{R}^{C}$, the
concatenation $[z; e_s]$ passes through a two-layer bottleneck MLP
with a sigmoid output $\alpha \in (0,1)^{C}$, and the module returns
$F' = F \odot \alpha$ with a residual connection. Conditioning the
gate computation on $e_s$ lets the network re-weight feature channels
differently for Urdu, English, and digit inputs, which is the
mechanism the script-statistics argument of the Introduction
predicts should matter, particularly for Urdu's dot- and
stroke-formation cues. Insertion depth is a swept factor, not a fixed
choice: modules are inserted at early, mid, late, or all three
boundaries, and the sweep is applied identically to the conditioned
arm and its matched control so that neither side is handicapped.

\subsection{Experimental arms}

\begin{description}
  \item[A0, re-run baselines.] The backbones and heads of
  \citet{kashif2026}, re-implemented and re-run under our fixed split
  and training protocol, anchoring all comparisons internally.
  \item[A1, matched-capacity script-agnostic control.] The backbone
  with SCA modules in which the script embedding $e_s$ is replaced by
  a single \emph{learned constant} embedding shared by all inputs.
  Parameter count, architecture, and tuning budget are identical to
  A2; only the script information is removed. This is the primary
  control.
  \item[A2, script-conditioned model (ours).] The backbone with SCA
  modules receiving the annotated (oracle) script tag.
  \item[A2$'$, predicted-script variant.] A2 with the oracle tag
  replaced by the output of a lightweight script classifier trained
  on the same split, quantifying robustness of any gain to imperfect
  script identification.
  \item[A3, naive conditioning.] The backbone without SCA modules,
  with the one-hot script tag concatenated to the penultimate
  features, in the spirit of simple feature-wise conditioning
  \citep{perez2018film}. Separates the value of the script
  \emph{information} from the value of the attention \emph{mechanism}
  that uses it.
  \item[A4, per-script experts.] Independent models trained on the
  Urdu-only and English-only subsets (digits held aside), routed by
  script at inference, testing whether a joint conditioned model
  outperforms full separation.
  \item[A5, identify-then-branch.] A two-stage pipeline in the
  multi-script recognition style \citep{multrenets2020}: a script
  classifier routes each image to a script-specific head over shared
  frozen features, testing whether joint conditioning beats explicit
  routing.
\end{description}

The primary matched pair is \textbf{A2 versus A1}: identical
capacity, identical tuning, identical split, differing only in
whether the gates see the script. The secondary pair is A2 versus A3.
All remaining comparisons are reported as context.

\subsection{Fair tuning protocol}

Both arms of every reported pair receive the same hyperparameter
search space, the same search budget, and the same model-selection
rule, with selection performed exclusively on the validation
partition. For the depth sweep, each arm reports its own
best-validated insertion configuration; the comparison is therefore
best-fair-configuration against best-fair-configuration rather than
an arbitrary fixed setting, per the pre-specified design goal of
giving a true effect its best honest chance without favoring either
side. The full search grid, budgets, and selected configurations are
released with the code. \todonote{Fix the concrete grid (learning
rate, weight decay, embedding dim $d$, bottleneck ratio, epochs,
early stopping) before implementation and freeze it in the repo
before any test-set evaluation.}

% ============================================================
\section{Experimental Setup}
\label{sec:setup}

\subsection{Split and its honest scope}

All experiments use a single fixed 80:10:10 train, validation, test
partition at the image level, stratified jointly by class and script
so that every partition preserves the class balance and script
composition of the release. The partition is generated once with a
recorded seed and reused unchanged by every arm; no arm ever sees a
different split. Because the release contains no writer identifiers
(Section~\ref{sec:data}), images from the same child may appear on
both sides of the partition, and the site confound is present on both
sides by construction. All absolute numbers in
Section~\ref{sec:results} are therefore \emph{within-dataset}
performance and are not generalization estimates; the quantity this
paper is designed to measure is the paired difference between arms
under identical conditions.

\subsection{Repeated runs and a power-motivated repeat count}

Every arm is trained under $S$ random seeds governing initialization
and data ordering, with the split held fixed. The repeat count is
chosen for sensitivity rather than convention: a pilot of five seeds
on the primary pair (A2, A1) estimates the between-seed standard
deviation of the paired accuracy difference, and $S$ is then set so
that a paired test has approximately 80\% power to detect a
pre-specified minimal effect of interest of 1.5 accuracy points,
subject to a compute ceiling consistent with a single-GPU
Colab/Kaggle budget \todonote{insert the pilot SD, the resulting
$S$, and the ceiling actually used; pre-register the minimal effect
of interest in the repo before the pilot}. Seeds are identical across
arms so that differences are paired seed-by-seed.

\subsection{Metrics and endpoints}

For every arm we report accuracy, precision, recall, $F_1$, and AUC
as mean $\pm$ standard deviation over seeds, pooled and broken out
per script (Urdu, English, digits). Two co-primary endpoints are
pre-specified: (i) the pooled A2$-$A1 accuracy difference and (ii)
the Urdu-subset A2$-$A1 accuracy difference, reflecting the
mechanism's prediction that gains concentrate in Urdu; a
Holm correction is applied across these two endpoints. All other
comparisons, including A2 versus A3, per-script breakdowns beyond the
Urdu endpoint, depth-sweep observations, and A2$'$ robustness, are
reported as exploratory and labeled as such. This pre-specification
protects the analysis in both directions: a real Urdu-specific effect
cannot be averaged away into a flat pooled number, and a spurious
effect cannot be harvested from an unplanned subgroup.

\subsection{Statistical analysis}

Each co-primary endpoint is tested with a paired $t$-test across the
$S$ seed-matched differences, accompanied by a Wilcoxon signed-rank
test as a distribution-free check, the effect size (Cohen's $d_z$),
and a 95\% confidence interval for the mean difference. With $n=852$
images and image-level evaluation, intervals will be wide, and we
report them as such; a null or negative result is reported with the
same prominence as a positive one. No test-set evaluation occurs
until the tuning grid, endpoints, and analysis plan are frozen in the
repository.

\subsection{Reproducibility}

We release code, the fixed split indices, the script annotations,
all configuration files, per-seed raw results, and the exact
environment specification. Training uses TensorFlow/Keras to remain
close to the source study's stack \todonote{confirm framework choice
at implementation; if PyTorch is preferred for the SCA modules,
state the deviation and re-run A0 in the same framework so the
anchor stays internal}. Every number in Section~\ref{sec:results}
must be traceable to a logged run artifact.

% ============================================================
\section{Results}
\label{sec:results}
\todonote{Pending. No numbers may be written here except from
executed experiments supplied by the authors.}

% ============================================================
\section{Limitations}
\label{sec:limitations}
\todonote{Pending. Must include: within-dataset only; site/label
confounds; construct validity of institutional labels; equal-
separability assumption behind the matched-pair logic; no
writer-level CV possible on released data.}

% ============================================================
\section{Future Work}
\label{sec:future}
\todonote{Pending. Parked list incl.\ multi-site collection with
writer IDs, participant-wise CV, cross-dataset testing incl.\ the
Kaggle English letter corpus, per-marker annotation, word/line-level
samples, clinical-grade labels.}

% ============================================================
\section*{Declarations}
\todonote{Required back matter: Data Availability, Ethics, CRediT
author contributions, Conflicts of Interest, Funding, AI-use
statement. Drafted at Stage 5.}

\paragraph{Resource availability (links to be filled at submission).}
Code, configuration files, the fixed split indices, the script
annotation file, and per-seed raw results are available at
\todonote{GitHub repository URL} and archived with a versioned DOI at
\todonote{Zenodo DOI}. The handwriting dataset itself is available
from its original source \citep{kashif2026} via the access route
published there; we do not redistribute the images.

\bibliography{ScriptAwareDyslexia_TBD_v1_2026-09}

\end{document}
